\section{Related work}\label{sec:related-work}

\paragraph{Cantor sets, pressure and fixed rule systems.}
The classical theory builds Cantor sets from a rule system given in advance: an iterated function system \cite{Hutchinson1981Fractals}, a Moran construction \cite{Moran1946AdditiveFunctions}, or a graph-directed Markov system \cite{MauldinUrbanskiGDMS}. It then relates their geometry to a pressure function \cite{Bowen1979Quasicircles,Ruelle1978ThermodynamicFormalism,FalconerFractalGeometry}. Random, non-autonomous and local versions let the maps vary with time, location or context \cite{FalconerRandomFractals1986,BarnsleyDemkoEltonGeronimo1988,RempeGillenUrbanski2016NonAutonomousCIFS,OliveiraVarandasLocalIFS}. In our setting the rule system itself, the kernel, is rewritten by the dynamics. The Cantor structure enters through the space of admissible futures, not through a limit set, and we make no dimension claim.

\paragraph{Products of nonnegative matrices and subadditive pressure.}
The growth pressure of Section~\ref{sec:pressure} is a pressure of products of nonnegative matrices along a dynamical system. Feng and Lau studied such pressure functions for products of nonnegative matrices \cite{FengLau2002MatrixPressure}. The subadditive and almost-additive thermodynamic formalism \cite{BarreiraNonadditiveTF,CaoFengHuang2008SubAdditive,Mummert2006AlmostAdditive} and the classical results on matrix products and subadditive processes \cite{FurstenbergKesten1960RandomMatrices,Oseledets1968MET,Kingman1968Subadditive} supply the general framework. Our existence proof uses only Fekete's lemma \cite{Fekete1923DistributionRoots} and uniform entry bounds on the shell. The reference laws of Proposition~\ref{prop:positive-gap} rest on Schreiber's identification of supremum growth rates of subadditive cocycles with invariant-measure growth rates \cite{Schreiber1998Subadditive}.

\paragraph{Factors, relative pressure and disintegration.}
Non-factorization through a base map is the standard way to say that one description strictly refines another, as for factor maps in symbolic dynamics \cite{LindMarcus1995SymbolicDynamics}. Disintegrating a measure over the fibers of a factor goes back to Rokhlin \cite{Rokhlin1949FundamentalIdeas}. Relative pressure and the relativized variational principle \cite{LedrappierWalters1977RelativisedVP,Walters1986RelativePressure} are the natural setting for comparing a global pressure with fiberwise ones. Our Section~\ref{sec:disintegration} works with the simplest instance, finitely many fibers of positive mass, where the comparison reduces to Proposition~\ref{prop:finite-disintegration}. The point there is not a new formula. It is that the fibers of a strict structural extension need not separate growth rates.

\paragraph{Lumpability and macro-closure.}
Whether coarse-grained masses evolve by a closed Markov dynamics is the classical lumpability question \cite{KemenySnell1976FiniteMarkovChains}. The completion map of Section~\ref{sec:completion} forgets the distribution within cells and reinstates it from prototypes. Its fixed point is defined whether or not the cells are lumpable, and at the warm states they are not (Remark~\ref{rem:refinement}).

\paragraph{The Six Birds program.}
The six mechanisms and the evolve--forget--reinstate completion come from the Six Birds program \cite{TsiokosFoundations,TsiokosCreateStone}. There, strict theory extension is defined as non-factorization with retention of the old description, and closure deficits are studied through predictive Kullback--Leibler losses \cite{TsiokosCastStone,CoverThomas2006InfoTheory}. Those results are structural or predictive statements about a specified finite host. They do not contain a pressure functional. The present paper supplies a concrete continuous-state instance of a strict retentive extension. It also shows that, in this instance, the extension is invisible to asymptotic growth. This complements those structural results and does not contradict them. For the warm pair under a common future law, the finite-horizon distinction, such as the nonzero two-step partition difference in Section~\ref{sec:zero-gap}, contributes only a bounded factor, because the history products coincide from the second step and the initial vectors have full support. It therefore disappears from exponential growth rates. In general, a predictive distinction can of course persist at a positive rate.

\paragraph{Hybrid and switched systems.}
Our substrate is a deterministic hybrid system with switching and hysteresis, driven by an input word \cite{GoebelSanfeliceTeel2012Hybrid}. The controlled shell is an invariant set for prescribed inputs, in the spirit of controlled invariance. It is not a statement about a random invariant set under an input distribution.
